% =====================================================================
% Pruebas funcionales de RF-02 (Facultades, carreras y usuarios) - Kamila
% FRAGMENTO para el informe: % =====================================================================
% Pruebas funcionales de RF-02 (Facultades, carreras y usuarios) - Kamila
% FRAGMENTO para el informe: % =====================================================================
% Pruebas funcionales de RF-02 (Facultades, carreras y usuarios) - Kamila
% FRAGMENTO para el informe: % =====================================================================
% Pruebas funcionales de RF-02 (Facultades, carreras y usuarios) - Kamila
% FRAGMENTO para el informe: \input{rf02_tabla_pruebas}
% Paquetes: booktabs, longtable, array. Las capturas estan en la carpeta evidencias/.
% =====================================================================

\providecommand{\evus}{\textunderscore\allowbreak}
\providecommand{\ev}[1]{{\texttt{\scriptsize\let\_\evus #1}}}

\subsection{Pruebas funcionales de RF-02: facultades, carreras y usuarios}

Cada prueba se hizo desde la aplicación y después se comprobó en PostgreSQL con una
consulta, para confirmar que lo que se vio en pantalla coincide con lo guardado en la base
de datos. Las capturas están en la carpeta \texttt{evidencias/} del repositorio.

\begingroup\footnotesize
\setlength{\tabcolsep}{3pt}
\begin{longtable}{>{\raggedright\arraybackslash}p{1.5cm}>{\raggedright\arraybackslash}p{1.8cm}>{\raggedright\arraybackslash}p{2.8cm}>{\raggedright\arraybackslash}p{2.7cm}>{\raggedright\arraybackslash}p{2.7cm}>{\raggedright\arraybackslash}p{3.2cm}}
\caption{Pruebas funcionales de RF-02}\label{tab:pruebas-rf02}\\
\toprule
\textbf{Requisito} & \textbf{Prueba} & \textbf{Datos usados} & \textbf{Resultado esperado} & \textbf{Resultado obtenido} & \textbf{Evidencia} \\
\midrule
\endfirsthead
\toprule
\textbf{Requisito} & \textbf{Prueba} & \textbf{Datos usados} & \textbf{Resultado esperado} & \textbf{Resultado obtenido} & \textbf{Evidencia} \\
\midrule
\endhead
RF-02 & Crear facultad &
Nombre: Facultad de Ciencias de la Ingeniería; activa. &
Se guarda y aparece en la lista con el mensaje ``Facultad creada''. &
Apareció el mensaje y la facultad en la lista. En PostgreSQL, \texttt{facultad} tiene la fila con \texttt{id\_facultad = 1}. &
\ev{RF02\_crear\_facultad.png}\newline
\ev{RF02\_crear\_facultad\_sql.png} \\
\addlinespace
RF-02 & Crear carrera &
Nombre: Ingeniería de Software; facultad elegida de la lista. &
Se guarda asociada a su facultad, con el mensaje ``Carrera creada''. &
Apareció el mensaje y la carrera junto a su facultad. En PostgreSQL, la consulta con \texttt{JOIN} muestra la carrera con su facultad. &
\ev{RF02\_crear\_carrera.png}\newline
\ev{RF02\_crear\_carrera\_sql.png} \\
\addlinespace
RF-02 & Crear usuario tutor &
Nombres: Tutor Prueba Uno; correo: tutor1@uteq.edu.ec; rol: tutor; carrera: Ingeniería de Software. &
Se guarda con su rol y carrera, con el mensaje ``Usuario creado''. &
Apareció el mensaje y el usuario en la lista. En PostgreSQL, \texttt{usuario} muestra el tutor con su rol y su carrera. &
\ev{RF02\_crear\_usuario\_tutor.png}\newline
\ev{RF02\_crear\_usuario\_tutor\_sql.png} \\
\addlinespace
RF-02 & Error: correo repetido &
Nombres: Tutor Prueba Uno; correo: tutor1@uteq.edu.ec (ya registrado); rol: tutor; misma carrera. &
No se guarda y se muestra ``Ya existe un usuario con ese correo''. &
El sistema rechazó el registro y mostró el mensaje bajo el campo Correo. En PostgreSQL, el correo aparece una sola vez. &
\ev{RF02\_error\_correo\_repetido.png}\newline
\ev{RF02\_error\_correo\_repetido\_sql.png} \\
\addlinespace
RF-02 & Error: eliminar facultad con carreras &
Facultad de Ciencias de la Ingeniería, que tiene 1 carrera asociada. &
No se elimina y se muestra ``No se puede eliminar: la facultad tiene carreras asociadas''. &
El sistema mostró el mensaje y la facultad siguió en la lista. En PostgreSQL, la facultad sigue con 1 carrera. &
\ev{RF02\_error\_eliminar\_facultad.png}\newline
\ev{RF02\_error\_eliminar\_facultad\_sql.png} \\
\addlinespace
RF-02 & Error: tutor sin carrera &
Nombres: Tutor Sin Carrera; correo: tutor3@uteq.edu.ec; rol: tutor; sin carrera. &
No se guarda y se muestra ``Solo el administrador general puede no tener carrera''. &
El sistema rechazó el registro y mostró el mensaje bajo el campo Carrera. &
\ev{RF02\_error\_tutor\_sin\_carrera.png} \\
\addlinespace
RF-02 & Error: eliminar usuario con documentos &
Tutor Prueba Uno, que tiene 1 documento asociado (Informe de prueba, PDF). &
No se elimina y se muestra ``No se puede eliminar: el usuario tiene documentos asociados''. &
El sistema mostró el mensaje y el usuario siguió en la lista. En PostgreSQL, el tutor sigue con 1 documento. &
\ev{RF02\_error\_eliminar\_usuario\_con\_documentos.png}\newline
\ev{RF02\_error\_eliminar\_usuario\_con\_documentos\_sql.png} \\
\bottomrule
\end{longtable}
\endgroup

Las consultas usadas para comprobar los datos en PostgreSQL fueron:

\begingroup\sloppy\raggedright
\begin{itemize}[leftmargin=1.2cm]
  \item \ev{SELECT * FROM facultad;}
  \item \ev{SELECT c.id\_carrera, c.nombre, f.nombre AS facultad FROM carrera c JOIN
        facultad f ON f.id\_facultad = c.id\_facultad;}
  \item \ev{SELECT u.id\_usuario, u.nombres, u.correo, r.nombre AS rol, c.nombre AS
        carrera FROM usuario u JOIN rol r ON r.id\_rol = u.id\_rol LEFT JOIN carrera c ON
        c.id\_carrera = u.id\_carrera;}
  \item \ev{SELECT correo, COUNT(*) AS veces FROM usuario GROUP BY correo;}
  \item \ev{SELECT f.id\_facultad, f.nombre, COUNT(c.id\_carrera) AS carreras FROM
        facultad f LEFT JOIN carrera c ON c.id\_facultad = f.id\_facultad GROUP BY
        f.id\_facultad, f.nombre;}
  \item \ev{SELECT u.id\_usuario, u.nombres, COUNT(d.id\_documento) AS documentos FROM
        usuario u LEFT JOIN documento d ON d.id\_usuario = u.id\_usuario GROUP BY
        u.id\_usuario, u.nombres;}
\end{itemize}
\endgroup

La prueba de eliminar un usuario que tiene documentos se hizo cuando el módulo de documentos (RF-03) ya estaba unido a la rama principal.

% Paquetes: booktabs, longtable, array. Las capturas estan en la carpeta evidencias/.
% =====================================================================

\providecommand{\evus}{\textunderscore\allowbreak}
\providecommand{\ev}[1]{{\texttt{\scriptsize\let\_\evus #1}}}

\subsection{Pruebas funcionales de RF-02: facultades, carreras y usuarios}

Cada prueba se hizo desde la aplicación y después se comprobó en PostgreSQL con una
consulta, para confirmar que lo que se vio en pantalla coincide con lo guardado en la base
de datos. Las capturas están en la carpeta \texttt{evidencias/} del repositorio.

\begingroup\footnotesize
\setlength{\tabcolsep}{3pt}
\begin{longtable}{>{\raggedright\arraybackslash}p{1.5cm}>{\raggedright\arraybackslash}p{1.8cm}>{\raggedright\arraybackslash}p{2.8cm}>{\raggedright\arraybackslash}p{2.7cm}>{\raggedright\arraybackslash}p{2.7cm}>{\raggedright\arraybackslash}p{3.2cm}}
\caption{Pruebas funcionales de RF-02}\label{tab:pruebas-rf02}\\
\toprule
\textbf{Requisito} & \textbf{Prueba} & \textbf{Datos usados} & \textbf{Resultado esperado} & \textbf{Resultado obtenido} & \textbf{Evidencia} \\
\midrule
\endfirsthead
\toprule
\textbf{Requisito} & \textbf{Prueba} & \textbf{Datos usados} & \textbf{Resultado esperado} & \textbf{Resultado obtenido} & \textbf{Evidencia} \\
\midrule
\endhead
RF-02 & Crear facultad &
Nombre: Facultad de Ciencias de la Ingeniería; activa. &
Se guarda y aparece en la lista con el mensaje ``Facultad creada''. &
Apareció el mensaje y la facultad en la lista. En PostgreSQL, \texttt{facultad} tiene la fila con \texttt{id\_facultad = 1}. &
\ev{RF02\_crear\_facultad.png}\newline
\ev{RF02\_crear\_facultad\_sql.png} \\
\addlinespace
RF-02 & Crear carrera &
Nombre: Ingeniería de Software; facultad elegida de la lista. &
Se guarda asociada a su facultad, con el mensaje ``Carrera creada''. &
Apareció el mensaje y la carrera junto a su facultad. En PostgreSQL, la consulta con \texttt{JOIN} muestra la carrera con su facultad. &
\ev{RF02\_crear\_carrera.png}\newline
\ev{RF02\_crear\_carrera\_sql.png} \\
\addlinespace
RF-02 & Crear usuario tutor &
Nombres: Tutor Prueba Uno; correo: tutor1@uteq.edu.ec; rol: tutor; carrera: Ingeniería de Software. &
Se guarda con su rol y carrera, con el mensaje ``Usuario creado''. &
Apareció el mensaje y el usuario en la lista. En PostgreSQL, \texttt{usuario} muestra el tutor con su rol y su carrera. &
\ev{RF02\_crear\_usuario\_tutor.png}\newline
\ev{RF02\_crear\_usuario\_tutor\_sql.png} \\
\addlinespace
RF-02 & Error: correo repetido &
Nombres: Tutor Prueba Uno; correo: tutor1@uteq.edu.ec (ya registrado); rol: tutor; misma carrera. &
No se guarda y se muestra ``Ya existe un usuario con ese correo''. &
El sistema rechazó el registro y mostró el mensaje bajo el campo Correo. En PostgreSQL, el correo aparece una sola vez. &
\ev{RF02\_error\_correo\_repetido.png}\newline
\ev{RF02\_error\_correo\_repetido\_sql.png} \\
\addlinespace
RF-02 & Error: eliminar facultad con carreras &
Facultad de Ciencias de la Ingeniería, que tiene 1 carrera asociada. &
No se elimina y se muestra ``No se puede eliminar: la facultad tiene carreras asociadas''. &
El sistema mostró el mensaje y la facultad siguió en la lista. En PostgreSQL, la facultad sigue con 1 carrera. &
\ev{RF02\_error\_eliminar\_facultad.png}\newline
\ev{RF02\_error\_eliminar\_facultad\_sql.png} \\
\addlinespace
RF-02 & Error: tutor sin carrera &
Nombres: Tutor Sin Carrera; correo: tutor3@uteq.edu.ec; rol: tutor; sin carrera. &
No se guarda y se muestra ``Solo el administrador general puede no tener carrera''. &
El sistema rechazó el registro y mostró el mensaje bajo el campo Carrera. &
\ev{RF02\_error\_tutor\_sin\_carrera.png} \\
\addlinespace
RF-02 & Error: eliminar usuario con documentos &
Tutor Prueba Uno, que tiene 1 documento asociado (Informe de prueba, PDF). &
No se elimina y se muestra ``No se puede eliminar: el usuario tiene documentos asociados''. &
El sistema mostró el mensaje y el usuario siguió en la lista. En PostgreSQL, el tutor sigue con 1 documento. &
\ev{RF02\_error\_eliminar\_usuario\_con\_documentos.png}\newline
\ev{RF02\_error\_eliminar\_usuario\_con\_documentos\_sql.png} \\
\bottomrule
\end{longtable}
\endgroup

Las consultas usadas para comprobar los datos en PostgreSQL fueron:

\begingroup\sloppy\raggedright
\begin{itemize}[leftmargin=1.2cm]
  \item \ev{SELECT * FROM facultad;}
  \item \ev{SELECT c.id\_carrera, c.nombre, f.nombre AS facultad FROM carrera c JOIN
        facultad f ON f.id\_facultad = c.id\_facultad;}
  \item \ev{SELECT u.id\_usuario, u.nombres, u.correo, r.nombre AS rol, c.nombre AS
        carrera FROM usuario u JOIN rol r ON r.id\_rol = u.id\_rol LEFT JOIN carrera c ON
        c.id\_carrera = u.id\_carrera;}
  \item \ev{SELECT correo, COUNT(*) AS veces FROM usuario GROUP BY correo;}
  \item \ev{SELECT f.id\_facultad, f.nombre, COUNT(c.id\_carrera) AS carreras FROM
        facultad f LEFT JOIN carrera c ON c.id\_facultad = f.id\_facultad GROUP BY
        f.id\_facultad, f.nombre;}
  \item \ev{SELECT u.id\_usuario, u.nombres, COUNT(d.id\_documento) AS documentos FROM
        usuario u LEFT JOIN documento d ON d.id\_usuario = u.id\_usuario GROUP BY
        u.id\_usuario, u.nombres;}
\end{itemize}
\endgroup

La prueba de eliminar un usuario que tiene documentos se hizo cuando el módulo de documentos (RF-03) ya estaba unido a la rama principal.

% Paquetes: booktabs, longtable, array. Las capturas estan en la carpeta evidencias/.
% =====================================================================

\providecommand{\evus}{\textunderscore\allowbreak}
\providecommand{\ev}[1]{{\texttt{\scriptsize\let\_\evus #1}}}

\subsection{Pruebas funcionales de RF-02: facultades, carreras y usuarios}

Cada prueba se hizo desde la aplicación y después se comprobó en PostgreSQL con una
consulta, para confirmar que lo que se vio en pantalla coincide con lo guardado en la base
de datos. Las capturas están en la carpeta \texttt{evidencias/} del repositorio.

\begingroup\footnotesize
\setlength{\tabcolsep}{3pt}
\begin{longtable}{>{\raggedright\arraybackslash}p{1.5cm}>{\raggedright\arraybackslash}p{1.8cm}>{\raggedright\arraybackslash}p{2.8cm}>{\raggedright\arraybackslash}p{2.7cm}>{\raggedright\arraybackslash}p{2.7cm}>{\raggedright\arraybackslash}p{3.2cm}}
\caption{Pruebas funcionales de RF-02}\label{tab:pruebas-rf02}\\
\toprule
\textbf{Requisito} & \textbf{Prueba} & \textbf{Datos usados} & \textbf{Resultado esperado} & \textbf{Resultado obtenido} & \textbf{Evidencia} \\
\midrule
\endfirsthead
\toprule
\textbf{Requisito} & \textbf{Prueba} & \textbf{Datos usados} & \textbf{Resultado esperado} & \textbf{Resultado obtenido} & \textbf{Evidencia} \\
\midrule
\endhead
RF-02 & Crear facultad &
Nombre: Facultad de Ciencias de la Ingeniería; activa. &
Se guarda y aparece en la lista con el mensaje ``Facultad creada''. &
Apareció el mensaje y la facultad en la lista. En PostgreSQL, \texttt{facultad} tiene la fila con \texttt{id\_facultad = 1}. &
\ev{RF02\_crear\_facultad.png}\newline
\ev{RF02\_crear\_facultad\_sql.png} \\
\addlinespace
RF-02 & Crear carrera &
Nombre: Ingeniería de Software; facultad elegida de la lista. &
Se guarda asociada a su facultad, con el mensaje ``Carrera creada''. &
Apareció el mensaje y la carrera junto a su facultad. En PostgreSQL, la consulta con \texttt{JOIN} muestra la carrera con su facultad. &
\ev{RF02\_crear\_carrera.png}\newline
\ev{RF02\_crear\_carrera\_sql.png} \\
\addlinespace
RF-02 & Crear usuario tutor &
Nombres: Tutor Prueba Uno; correo: tutor1@uteq.edu.ec; rol: tutor; carrera: Ingeniería de Software. &
Se guarda con su rol y carrera, con el mensaje ``Usuario creado''. &
Apareció el mensaje y el usuario en la lista. En PostgreSQL, \texttt{usuario} muestra el tutor con su rol y su carrera. &
\ev{RF02\_crear\_usuario\_tutor.png}\newline
\ev{RF02\_crear\_usuario\_tutor\_sql.png} \\
\addlinespace
RF-02 & Error: correo repetido &
Nombres: Tutor Prueba Uno; correo: tutor1@uteq.edu.ec (ya registrado); rol: tutor; misma carrera. &
No se guarda y se muestra ``Ya existe un usuario con ese correo''. &
El sistema rechazó el registro y mostró el mensaje bajo el campo Correo. En PostgreSQL, el correo aparece una sola vez. &
\ev{RF02\_error\_correo\_repetido.png}\newline
\ev{RF02\_error\_correo\_repetido\_sql.png} \\
\addlinespace
RF-02 & Error: eliminar facultad con carreras &
Facultad de Ciencias de la Ingeniería, que tiene 1 carrera asociada. &
No se elimina y se muestra ``No se puede eliminar: la facultad tiene carreras asociadas''. &
El sistema mostró el mensaje y la facultad siguió en la lista. En PostgreSQL, la facultad sigue con 1 carrera. &
\ev{RF02\_error\_eliminar\_facultad.png}\newline
\ev{RF02\_error\_eliminar\_facultad\_sql.png} \\
\addlinespace
RF-02 & Error: tutor sin carrera &
Nombres: Tutor Sin Carrera; correo: tutor3@uteq.edu.ec; rol: tutor; sin carrera. &
No se guarda y se muestra ``Solo el administrador general puede no tener carrera''. &
El sistema rechazó el registro y mostró el mensaje bajo el campo Carrera. &
\ev{RF02\_error\_tutor\_sin\_carrera.png} \\
\addlinespace
RF-02 & Error: eliminar usuario con documentos &
Tutor Prueba Uno, que tiene 1 documento asociado (Informe de prueba, PDF). &
No se elimina y se muestra ``No se puede eliminar: el usuario tiene documentos asociados''. &
El sistema mostró el mensaje y el usuario siguió en la lista. En PostgreSQL, el tutor sigue con 1 documento. &
\ev{RF02\_error\_eliminar\_usuario\_con\_documentos.png}\newline
\ev{RF02\_error\_eliminar\_usuario\_con\_documentos\_sql.png} \\
\bottomrule
\end{longtable}
\endgroup

Las consultas usadas para comprobar los datos en PostgreSQL fueron:

\begingroup\sloppy\raggedright
\begin{itemize}[leftmargin=1.2cm]
  \item \ev{SELECT * FROM facultad;}
  \item \ev{SELECT c.id\_carrera, c.nombre, f.nombre AS facultad FROM carrera c JOIN
        facultad f ON f.id\_facultad = c.id\_facultad;}
  \item \ev{SELECT u.id\_usuario, u.nombres, u.correo, r.nombre AS rol, c.nombre AS
        carrera FROM usuario u JOIN rol r ON r.id\_rol = u.id\_rol LEFT JOIN carrera c ON
        c.id\_carrera = u.id\_carrera;}
  \item \ev{SELECT correo, COUNT(*) AS veces FROM usuario GROUP BY correo;}
  \item \ev{SELECT f.id\_facultad, f.nombre, COUNT(c.id\_carrera) AS carreras FROM
        facultad f LEFT JOIN carrera c ON c.id\_facultad = f.id\_facultad GROUP BY
        f.id\_facultad, f.nombre;}
  \item \ev{SELECT u.id\_usuario, u.nombres, COUNT(d.id\_documento) AS documentos FROM
        usuario u LEFT JOIN documento d ON d.id\_usuario = u.id\_usuario GROUP BY
        u.id\_usuario, u.nombres;}
\end{itemize}
\endgroup

La prueba de eliminar un usuario que tiene documentos se hizo cuando el módulo de documentos (RF-03) ya estaba unido a la rama principal.

% Paquetes: booktabs, longtable, array. Las capturas estan en la carpeta evidencias/.
% =====================================================================

\providecommand{\evus}{\textunderscore\allowbreak}
\providecommand{\ev}[1]{{\texttt{\scriptsize\let\_\evus #1}}}

\subsection{Pruebas funcionales de RF-02: facultades, carreras y usuarios}

Cada prueba se hizo desde la aplicación y después se comprobó en PostgreSQL con una
consulta, para confirmar que lo que se vio en pantalla coincide con lo guardado en la base
de datos. Las capturas están en la carpeta \texttt{evidencias/} del repositorio.

\begingroup\footnotesize
\setlength{\tabcolsep}{3pt}
\begin{longtable}{>{\raggedright\arraybackslash}p{1.5cm}>{\raggedright\arraybackslash}p{1.8cm}>{\raggedright\arraybackslash}p{2.8cm}>{\raggedright\arraybackslash}p{2.7cm}>{\raggedright\arraybackslash}p{2.7cm}>{\raggedright\arraybackslash}p{3.2cm}}
\caption{Pruebas funcionales de RF-02}\label{tab:pruebas-rf02}\\
\toprule
\textbf{Requisito} & \textbf{Prueba} & \textbf{Datos usados} & \textbf{Resultado esperado} & \textbf{Resultado obtenido} & \textbf{Evidencia} \\
\midrule
\endfirsthead
\toprule
\textbf{Requisito} & \textbf{Prueba} & \textbf{Datos usados} & \textbf{Resultado esperado} & \textbf{Resultado obtenido} & \textbf{Evidencia} \\
\midrule
\endhead
RF-02 & Crear facultad &
Nombre: Facultad de Ciencias de la Ingeniería; activa. &
Se guarda y aparece en la lista con el mensaje ``Facultad creada''. &
Apareció el mensaje y la facultad en la lista. En PostgreSQL, \texttt{facultad} tiene la fila con \texttt{id\_facultad = 1}. &
\ev{RF02\_crear\_facultad.png}\newline
\ev{RF02\_crear\_facultad\_sql.png} \\
\addlinespace
RF-02 & Crear carrera &
Nombre: Ingeniería de Software; facultad elegida de la lista. &
Se guarda asociada a su facultad, con el mensaje ``Carrera creada''. &
Apareció el mensaje y la carrera junto a su facultad. En PostgreSQL, la consulta con \texttt{JOIN} muestra la carrera con su facultad. &
\ev{RF02\_crear\_carrera.png}\newline
\ev{RF02\_crear\_carrera\_sql.png} \\
\addlinespace
RF-02 & Crear usuario tutor &
Nombres: Tutor Prueba Uno; correo: tutor1@uteq.edu.ec; rol: tutor; carrera: Ingeniería de Software. &
Se guarda con su rol y carrera, con el mensaje ``Usuario creado''. &
Apareció el mensaje y el usuario en la lista. En PostgreSQL, \texttt{usuario} muestra el tutor con su rol y su carrera. &
\ev{RF02\_crear\_usuario\_tutor.png}\newline
\ev{RF02\_crear\_usuario\_tutor\_sql.png} \\
\addlinespace
RF-02 & Error: correo repetido &
Nombres: Tutor Prueba Uno; correo: tutor1@uteq.edu.ec (ya registrado); rol: tutor; misma carrera. &
No se guarda y se muestra ``Ya existe un usuario con ese correo''. &
El sistema rechazó el registro y mostró el mensaje bajo el campo Correo. En PostgreSQL, el correo aparece una sola vez. &
\ev{RF02\_error\_correo\_repetido.png}\newline
\ev{RF02\_error\_correo\_repetido\_sql.png} \\
\addlinespace
RF-02 & Error: eliminar facultad con carreras &
Facultad de Ciencias de la Ingeniería, que tiene 1 carrera asociada. &
No se elimina y se muestra ``No se puede eliminar: la facultad tiene carreras asociadas''. &
El sistema mostró el mensaje y la facultad siguió en la lista. En PostgreSQL, la facultad sigue con 1 carrera. &
\ev{RF02\_error\_eliminar\_facultad.png}\newline
\ev{RF02\_error\_eliminar\_facultad\_sql.png} \\
\addlinespace
RF-02 & Error: tutor sin carrera &
Nombres: Tutor Sin Carrera; correo: tutor3@uteq.edu.ec; rol: tutor; sin carrera. &
No se guarda y se muestra ``Solo el administrador general puede no tener carrera''. &
El sistema rechazó el registro y mostró el mensaje bajo el campo Carrera. &
\ev{RF02\_error\_tutor\_sin\_carrera.png} \\
\addlinespace
RF-02 & Error: eliminar usuario con documentos &
Tutor Prueba Uno, que tiene 1 documento asociado (Informe de prueba, PDF). &
No se elimina y se muestra ``No se puede eliminar: el usuario tiene documentos asociados''. &
El sistema mostró el mensaje y el usuario siguió en la lista. En PostgreSQL, el tutor sigue con 1 documento. &
\ev{RF02\_error\_eliminar\_usuario\_con\_documentos.png}\newline
\ev{RF02\_error\_eliminar\_usuario\_con\_documentos\_sql.png} \\
\bottomrule
\end{longtable}
\endgroup

Las consultas usadas para comprobar los datos en PostgreSQL fueron:

\begingroup\sloppy\raggedright
\begin{itemize}[leftmargin=1.2cm]
  \item \ev{SELECT * FROM facultad;}
  \item \ev{SELECT c.id\_carrera, c.nombre, f.nombre AS facultad FROM carrera c JOIN
        facultad f ON f.id\_facultad = c.id\_facultad;}
  \item \ev{SELECT u.id\_usuario, u.nombres, u.correo, r.nombre AS rol, c.nombre AS
        carrera FROM usuario u JOIN rol r ON r.id\_rol = u.id\_rol LEFT JOIN carrera c ON
        c.id\_carrera = u.id\_carrera;}
  \item \ev{SELECT correo, COUNT(*) AS veces FROM usuario GROUP BY correo;}
  \item \ev{SELECT f.id\_facultad, f.nombre, COUNT(c.id\_carrera) AS carreras FROM
        facultad f LEFT JOIN carrera c ON c.id\_facultad = f.id\_facultad GROUP BY
        f.id\_facultad, f.nombre;}
  \item \ev{SELECT u.id\_usuario, u.nombres, COUNT(d.id\_documento) AS documentos FROM
        usuario u LEFT JOIN documento d ON d.id\_usuario = u.id\_usuario GROUP BY
        u.id\_usuario, u.nombres;}
\end{itemize}
\endgroup

La prueba de eliminar un usuario que tiene documentos se hizo cuando el módulo de documentos (RF-03) ya estaba unido a la rama principal.
